% Generated draft: requires mathematical and source review.
\documentclass[11pt]{article}
\usepackage[margin=25mm]{geometry}
\usepackage{fontspec}
\setmainfont{Latin Modern Roman}
\usepackage{amsmath,amssymb,mathtools,unicode-math}
\setmathfont{Latin Modern Math}
\usepackage{enumitem,needspace,xurl,hyperref}
\setlength{\parindent}{0pt}
\setlength{\parskip}{5pt}
\setlength{\emergencystretch}{3em}
\newcommand{\N}{\mathbb{N}}
\newcommand{\Z}{\mathbb{Z}}
\newcommand{\Q}{\mathbb{Q}}
\newcommand{\R}{\mathbb{R}}
\newcommand{\C}{\mathbb{C}}
\newcommand{\F}{\mathbb{F}}
\newcommand{\A}{\mathbb{A}}
\newcommand{\PP}{\mathbb{P}}
\newcommand{\E}{\mathbb{E}}
\newcommand{\eps}{\varepsilon}
\newcommand{\dd}{\,\mathrm{d}}
\begin{document}
\section*{UnsolvedMath ID: 800892}
\textbf{Source problem code:} catalogue entry
\subsection{Definitions and mathematical statement}
Let \(n\geq 1\) be an integer and let \(\Omega\subset\C^n\) be a bounded domain with smooth, strictly pseudoconvex boundary. Strict pseudoconvexity means that a smooth local defining function has positive Levi form on every nonzero complex tangent vector at the boundary. A function is plurisubharmonic if it is upper semicontinuous and its restriction to every complex line is subharmonic wherever defined.

Let \(u\) be plurisubharmonic on \(\Omega\), and assume
\[
u\in\mathcal C^{1,\alpha}(\Omega),\qquad \alpha>1-\frac{2}{n}.
\]
In the usual range \(0<\alpha\leq 1\), the Hölder condition means that the first derivatives are locally Hölder continuous with exponent \(\alpha\). The source does not explicitly specify the admissible range of \(\alpha\), or whether its Hölder assumption requires a finite global norm on \(\Omega\).

Suppose that there exist \(f\in\mathcal C^\infty(\Omega)\) and a constant \(\varepsilon>0\) such that
\[
(dd^c u)^n=f\,dV\quad\text{on }\Omega,
\qquad f(z)\geq\varepsilon\quad\text{for every }z\in\Omega,
\]
where \(dV\) is Euclidean volume measure. Since \(u\) is locally bounded and plurisubharmonic, the left side is interpreted as its complex Monge--Ampère measure. The normalization of \(d^c\) is not specified in the supplied source and must be used consistently with the density \(f\).

\textbf{Question [Blocki-Dinew].} Under these assumptions, must
\[
u\in\mathcal C^\infty(\Omega)?
\]
The conclusion concerns interior smoothness. No boundary values, extension to \(\overline\Omega\), or smoothness up to the boundary are specified.
\subsection{Short English statement}
Must a plurisubharmonic solution of the complex Monge--Ampère equation on a smooth bounded strictly pseudoconvex domain be smooth in the interior if its density is smooth and uniformly positive and its first derivatives have Hölder exponent \(\alpha>1-2/n\)? The source leaves the precise Hölder-space convention unspecified.
\subsection{Sources}
\begin{itemize}
\item Supplied local UnsolvedMath registry record: permanent ID \(800892\), source problem code \texttt{catalogue entry}, entitled \emph{Open Problems in Pluripotential Theory --- Question 23}. The record lists Dinew, Guedj, and Zeriahi as proposers, with proposed year 2015; the question itself carries the attribution \textup{[Blocki-Dinew]}.
\item Dinew, Guedj, and Zeriahi, \emph{Open Problems in Pluripotential Theory} (2015), Question 23. Supplied source URL: \url{https://arxiv.org/abs/1511.00705}. The supplied provenance records access on 2026-10-16 and extraction by \texttt{arxiv\_latex}; these are source-record claims, not independent access or verification performed for this draft.
\item The registry labels this problem \texttt{open}. This is supplied status information, not an independently verified assessment of the literature.
\end{itemize}
\end{document}
